% -*- coding=utf-8 -*-
\chapter{附录}

\section*{一、处理碰撞事例的Python代码}
%\begin{lstlisting}[language={python}, basicstyle = \ttfamily,breaklines = true, caption={Python 源代码}, label=code_1]
\begin{jllisting}[language={python}, caption={一段python源代码}, label=code:2]
kk=2;
[mdd,ndd]=size(dd);
while ~isempty(V)
    [tmpd,j]=min(W(i,V));tmpj=V(j);
for k=2:ndd
    [tmp1,jj]=min(dd(1,k)+W(dd(2,k),V));
    tmp2=V(jj);tt(k-1,:)=[tmp1,tmp2,jj];
end
    tmp=[tmpd,tmpj,j;tt];[tmp3,tmp4]=min(tmp(:,1));
if tmp3==tmpd, ss(1:2,kk)=[i;tmp(tmp4,2)];
 \end{jllisting}
 
 \newpage
 
 \section*{二、处理碰撞事例的Julia代码}
 \begin{jllisting}[language={julia}, caption={一段Julia源代码}, label=code:3]
  using LinearAlgebra
using Random

"""
PSO求解最小特征值及特征向量
A: 对称矩阵
n_particles: 粒子数量
max_iter: 最大迭代次数
"""
function pso_min_eigen(A::Matrix{Float64}; n_particles=50, max_iter=200)
    n = size(A, 1)

    # 定义目标函数：瑞利商 f(x) = (x'Ax) / (x'x)
    function rayleigh_quotient(x)
        return (x' * A * x) / (x' * x)
    end

    # 初始化粒子位置 (随机单位向量方向) 和速度
    particles = [randn(n) for _ in 1:n_particles]
    velocities = [zeros(n) for _ in 1:n_particles]

    # 个体最优位置及其适应度
    pbest_pos = copy(particles)
    pbest_val = [rayleigh_quotient(p) for p in particles]

    # 全局最优
    gbest_idx = argmin(pbest_val)
    gbest_pos = copy(pbest_pos[gbest_idx])
    gbest_val = pbest_val[gbest_idx]

    # 超参数
    w = 0.5   # 惯性权重
    c1 = 1.5  # 个体学习因子
    c2 = 1.5  # 全局学习因子

    for _ in 1:max_iter
        for i in 1:n_particles
            # 更新速度
            velocities[i] = w * velocities[i] + 
                            c1 * rand() * (pbest_pos[i] - particles[i]) + 
                            c2 * rand() * (gbest_pos - particles[i])

            # 更新位置
            particles[i] += velocities[i]

            # 计算适应度
            val = rayleigh_quotient(particles[i])

            # 更新个体最优
            if val < pbest_val[i]
                pbest_val[i] = val
                pbest_pos[i] = copy(particles[i])

                # 更新全局最优
                if val < gbest_val
                    gbest_val = val
                    gbest_pos = copy(particles[i])
                end
            end
        end
    end

    # 返回结果：特征值与单位特征向量
    return gbest_val, gbest_pos / norm(gbest_pos)
end

# --- 测试 ---
Random.seed!(123)
A = rand(5, 5)
A = A' * A # 构造对称正定矩阵
λ, v = pso_min_eigen(A)

println("PSO计算的最小特征值: ", λ)
println("对比线性代数库计算: ", minimum(eigvals(A)))
println("单位特征向量: ", v)
\end{jllisting}


\begin{jllisting}[language={julia}, caption={PSO求解矩阵特征值及其单位特征向量}, label=code:grok]
using LinearAlgebra
using Random

"""
使用 PSO（粒子群优化）求解矩阵 A 的**最小特征值**及其对应的**单位特征向量**。

实现思路：
- 在单位球面上优化目标函数 f(x) = x' * A * x（即 Rayleigh 商在 ||x||=1 时的形式）
- 粒子位置始终归一化到单位向量
"""

function pso_min_eigenvalue(A::AbstractMatrix{T};
                            n_particles::Int = 50,
                            max_iter::Int = 1000,
                            w::Float64 = 0.729,
                            c1::Float64 = 1.49445,
                            c2::Float64 = 1.49445,
                            w_decay::Float64 = 0.995,
                            seed::Union{Int, Nothing} = 42) where T <: Real

    n = size(A, 1)
    if seed !== nothing
        Random.seed!(seed)
    end

    # ==================== 初始化 ====================
    # 随机初始化粒子位置，并归一化到单位球面
    pos = randn(n, n_particles)
    pos .= pos ./ mapslices(norm, pos; dims = 1)

    # 初始化速度（较小）
    vel = randn(n, n_particles) .* 0.1

    # 个人最优
    pbest_pos = copy(pos)
    pbest_val = [dot(pos[:, i], A * pos[:, i]) for i in 1:n_particles]

    # 全局最优
    gbest_idx = argmin(pbest_val)
    gbest_pos = copy(pbest_pos[:, gbest_idx])
    gbest_val = pbest_val[gbest_idx]

    current_w = w

    # ==================== PSO 主循环 ====================
    for iter in 1:max_iter
        # 1. 更新个体最优和全局最优
        for i in 1:n_particles
            x = @view pos[:, i]
            val = dot(x, A * x)

            if val < pbest_val[i]
                pbest_val[i] = val
                pbest_pos[:, i] .= x
            end

            if val < gbest_val
                gbest_val = val
                gbest_pos .= x
            end
        end

        # 2. 更新速度和位置
        for i in 1:n_particles
            r1 = rand()
            r2 = rand()

            vel[:, i] .= current_w .* vel[:, i] .+
                         c1 * r1 .* (pbest_pos[:, i] .- pos[:, i]) .+
                         c2 * r2 .* (gbest_pos .- pos[:, i])

            pos[:, i] .+= vel[:, i]

            # 重新归一化到单位向量
            nrm = norm(pos[:, i])
            if nrm > 1e-12
                pos[:, i] ./= nrm
            end
        end

        # 惯性权重衰减
        current_w *= w_decay
    end

    # ==================== 返回结果 ====================
    unit_vec = gbest_pos ./ norm(gbest_pos)
    eigenvalue = dot(unit_vec, A * unit_vec)

    return eigenvalue, unit_vec
end


# ==================== 使用示例 ====================
A = Float64[4 2 1; 2 5 3; 1 3 6]

println("矩阵 A = ")
display(A)

λ, v = pso_min_eigenvalue(A, n_particles=60, max_iter=1200)

println("\nPSO 求得的最小特征值 ≈ ", λ)
println("对应的单位特征向量 ≈ ", v)
println("特征向量模长 = ", norm(v))
println("\n验证 A*v ≈ λ*v : ", isapprox(A*v, λ*v, atol=1e-3))
println("验证 Rayleigh 商: ", dot(v, A*v))

# 可选：与真实特征值对比
println("\n真实最小特征值（Julia 内置）: ", minimum(eigvals(A)))
\end{jllisting}

\begin{jllisting}[language={julia}, caption={PSO求解矩阵特征值及其单位特征向量}, label=code:grok1]
function pso_min_eigenvalue_multiple(A; n_runs=10, kwargs...)
    best_val = Inf
    best_vec = nothing
    for _ in 1:n_runs
        λ, v = pso_min_eigenvalue(A; kwargs...)
        if λ < best_val
            best_val = λ
            best_vec = v
        end
    end
    return best_val, best_vec
end
\end{jllisting}